% Options for packages loaded elsewhere
\PassOptionsToPackage{unicode}{hyperref}
\PassOptionsToPackage{hyphens}{url}
\PassOptionsToPackage{dvipsnames,svgnames*,x11names*}{xcolor}
%
\documentclass[
  10pt,
]{article}
\usepackage{lmodern}
\usepackage{amssymb,amsmath}
\usepackage{ifxetex,ifluatex}
\ifnum 0\ifxetex 1\fi\ifluatex 1\fi=0 % if pdftex
  \usepackage[T1]{fontenc}
  \usepackage[utf8]{inputenc}
  \usepackage{textcomp} % provide euro and other symbols
\else % if luatex or xetex
  \usepackage{unicode-math}
  \defaultfontfeatures{Scale=MatchLowercase}
  \defaultfontfeatures[\rmfamily]{Ligatures=TeX,Scale=1}
  \setmainfont[]{DejaVu Serif}
  \setmonofont[]{DejaVu Sans Mono}
\fi
% Use upquote if available, for straight quotes in verbatim environments
\IfFileExists{upquote.sty}{\usepackage{upquote}}{}
\IfFileExists{microtype.sty}{% use microtype if available
  \usepackage[]{microtype}
  \UseMicrotypeSet[protrusion]{basicmath} % disable protrusion for tt fonts
}{}
\makeatletter
\@ifundefined{KOMAClassName}{% if non-KOMA class
  \IfFileExists{parskip.sty}{%
    \usepackage{parskip}
  }{% else
    \setlength{\parindent}{0pt}
    \setlength{\parskip}{6pt plus 2pt minus 1pt}}
}{% if KOMA class
  \KOMAoptions{parskip=half}}
\makeatother
\usepackage{xcolor}
\IfFileExists{xurl.sty}{\usepackage{xurl}}{} % add URL line breaks if available
\IfFileExists{bookmark.sty}{\usepackage{bookmark}}{\usepackage{hyperref}}
\hypersetup{
  colorlinks=true,
  linkcolor=Maroon,
  filecolor=Maroon,
  citecolor=Blue,
  urlcolor=Blue,
  pdfcreator={LaTeX via pandoc}}
\urlstyle{same} % disable monospaced font for URLs
\usepackage[margin=0.6in,landscape]{geometry}
\usepackage{longtable,booktabs}
% Correct order of tables after \paragraph or \subparagraph
\usepackage{etoolbox}
\makeatletter
\patchcmd\longtable{\par}{\if@noskipsec\mbox{}\fi\par}{}{}
\makeatother
% Allow footnotes in longtable head/foot
\IfFileExists{footnotehyper.sty}{\usepackage{footnotehyper}}{\usepackage{footnote}}
\makesavenoteenv{longtable}
\setlength{\emergencystretch}{3em} % prevent overfull lines
\providecommand{\tightlist}{%
  \setlength{\itemsep}{0pt}\setlength{\parskip}{0pt}}
\setcounter{secnumdepth}{-\maxdimen} % remove section numbering
\renewcommand{\arraystretch}{1.6}

\author{}
\date{}

\begin{document}

\hypertarget{stage-1.2-non-constant-isp-ascent-hard-mode}{%
\section{Stage 1.2 --- Non-Constant Isp Ascent (Hard
Mode)}\label{stage-1.2-non-constant-isp-ascent-hard-mode}}

\textbf{Era:} Redstone/Mercury-Redstone class, \textasciitilde1961
(fictional program), continuing Project Arrowhead. \textbf{Target time:}
a full afternoon. \textbf{Prerequisite:} Stage 1.0 (Range Table
Targeting), Stage 1.1 (Liftoff Check \& Gravity-Turn Loss).

\hypertarget{why-this-problem-exists}{%
\subsection{Why this problem exists}\label{why-this-problem-exists}}

1.0 treated specific impulse as a single constant, 235 s, for the whole
burn. Real engines don't work that way: thrust is F = ṁ·Ve + (Pe −
Pamb)·Ae --- as the vehicle climbs and ambient pressure Pamb drops, the
same engine produces more thrust and a higher effective Isp, right up to
a vacuum ceiling. This is standard, period-authentic propulsion analysis
(see Sutton, \emph{Rocket Propulsion Elements} --- in print and in
active engineering use well before 1961) and it is \textbf{not} a small
effect: most of the pressure drop happens in the first
\textasciitilde100,000--150,000 ft, which is most of this vehicle's
burn. Treating Isp as constant, as 1.0 did, understates burnout velocity
by a meaningful margin --- you'll see by how much below.

This is genuinely harder than 1.0, for a specific reason: Isp now
depends on altitude, altitude depends on the velocity history, and
velocity depends on Isp. That circular dependency is why this can't be
solved in one closed-form step --- it has to be marched forward in time,
a handful of steps at a time, the same way early trajectory integrators
(Gauss-Jackson and similar --- already flagged in the campaign's
reference list as a harder future problem in its own right) worked
before continuous integration was practical by hand. What you're doing
here is a simplified version of that same idea.

\hypertarget{given-data}{%
\subsection{Given data}\label{given-data}}

Same (corrected) vehicle as 1.0/1.1: W₀ = 55,000 lbf, Wₚ = 40,000 lbf
burned uniformly over t\_b = 143 s (weight decreases linearly with
time). Standard gravity g₀ = 32.174 ft/s².

\textbf{Loss budget:} Δv\_losses = 3540 ft/s, from Stage 1.1's
first-principles gravity-turn/drag study.

\textbf{Engine performance --- Isp vs.~altitude} (from separate
propulsion analysis; reflects the nozzle's altitude compensation as
ambient pressure drops):

\begin{longtable}[]{@{}ll@{}}
\toprule
h (ft) & Isp (s)\tabularnewline
\midrule
\endhead
0 & 235\tabularnewline
20,000 & 246\tabularnewline
40,000 & 253\tabularnewline
60,000 & 258\tabularnewline
80,000 & 261\tabularnewline
100,000 & 263\tabularnewline
120,000 & 264\tabularnewline
≥ 140,000 & 265 (vacuum plateau)\tabularnewline
\bottomrule
\end{longtable}

\textbf{Ascent altitude profile:} reuse the same trajectory checkpoints
from Stage 1.1 (same vehicle, same preliminary shaping run) rather than
a separately-invented profile:

\begin{longtable}[]{@{}ll@{}}
\toprule
t (s) & h (ft)\tabularnewline
\midrule
\endhead
0 & 0\tabularnewline
13 & 604\tabularnewline
39 & 6,485\tabularnewline
65 & 19,230\tabularnewline
91 & 37,878\tabularnewline
117 & 63,427\tabularnewline
143 & 100,235\tabularnewline
\bottomrule
\end{longtable}

\textbf{Explicit simplification, stated up front:} this problem does
\emph{not} solve for the trajectory shape --- the altitude profile above
is a given input (1.1's result), not something you derive from the
velocity you're computing here. That keeps the coupling one-directional
(altitude → Isp → Δv) and hand-tractable. The full version, where
Isp-vs-altitude feeds back into the trajectory shape itself, is real and
harder --- a natural candidate for a future stage, not this one.

\hypertarget{authentic-method-reference-material}{%
\subsection{Authentic method + reference
material}\label{authentic-method-reference-material}}

Divide the burn into the six intervals defined by the checkpoints above
(note: unequal step sizes --- the first is 13 s, the rest are 26 s each,
since they come from 1.1's checkpoint times, not an even split). For
each step:

\begin{verbatim}
1. Look up altitude at the START of the step (from the given profile).
2. Interpolate Isp at that altitude.
3. c = Isp · g₀
4. Δv_step = c · ln(W_start/W_end)     [W_end from the linear weight burn]
5. Accumulate: V_cumulative += Δv_step
\end{verbatim}

\textbf{Note the numerical choice being made in step 1:} using the
altitude at the \emph{start} of each step to look up Isp for that whole
step is an explicit (forward-lagged) scheme --- it's simple and
hand-tractable, but it means each step's Isp is slightly stale (the
vehicle is actually higher, and Isp actually slightly better, by the
step's end). A more accurate scheme would use the midpoint or average
altitude of the step. This is the same kind of approximation-choice
tradeoff numerical integrators always face --- flagged here so it's a
visible decision, not a hidden one.

\hypertarget{worksheet}{%
\subsection{Worksheet}\label{worksheet}}

\begin{longtable}[]{@{}llllllllll@{}}
\toprule
\begin{minipage}[b]{0.03\columnwidth}\raggedright
Step\strut
\end{minipage} & \begin{minipage}[b]{0.10\columnwidth}\raggedright
t\_start--t\_end (s)\strut
\end{minipage} & \begin{minipage}[b]{0.07\columnwidth}\raggedright
h\_start (ft)\strut
\end{minipage} & \begin{minipage}[b]{0.08\columnwidth}\raggedright
Isp (interp.)\strut
\end{minipage} & \begin{minipage}[b]{0.06\columnwidth}\raggedright
c = Isp·g₀\strut
\end{minipage} & \begin{minipage}[b]{0.05\columnwidth}\raggedright
W\_start\strut
\end{minipage} & \begin{minipage}[b]{0.04\columnwidth}\raggedright
W\_end\strut
\end{minipage} & \begin{minipage}[b]{0.10\columnwidth}\raggedright
ln(W\_start/W\_end)\strut
\end{minipage} & \begin{minipage}[b]{0.09\columnwidth}\raggedright
Δv\_step (ft/s)\strut
\end{minipage} & \begin{minipage}[b]{0.11\columnwidth}\raggedright
V\_cumulative (ft/s)\strut
\end{minipage}\tabularnewline
\midrule
\endhead
\begin{minipage}[t]{0.03\columnwidth}\raggedright
1\strut
\end{minipage} & \begin{minipage}[t]{0.10\columnwidth}\raggedright
0--13\strut
\end{minipage} & \begin{minipage}[t]{0.07\columnwidth}\raggedright
0\strut
\end{minipage} & \begin{minipage}[t]{0.08\columnwidth}\raggedright
\strut
\end{minipage} & \begin{minipage}[t]{0.06\columnwidth}\raggedright
\strut
\end{minipage} & \begin{minipage}[t]{0.05\columnwidth}\raggedright
55,000\strut
\end{minipage} & \begin{minipage}[t]{0.04\columnwidth}\raggedright
51,364\strut
\end{minipage} & \begin{minipage}[t]{0.10\columnwidth}\raggedright
\strut
\end{minipage} & \begin{minipage}[t]{0.09\columnwidth}\raggedright
\strut
\end{minipage} & \begin{minipage}[t]{0.11\columnwidth}\raggedright
\strut
\end{minipage}\tabularnewline
\begin{minipage}[t]{0.03\columnwidth}\raggedright
2\strut
\end{minipage} & \begin{minipage}[t]{0.10\columnwidth}\raggedright
13--39\strut
\end{minipage} & \begin{minipage}[t]{0.07\columnwidth}\raggedright
604\strut
\end{minipage} & \begin{minipage}[t]{0.08\columnwidth}\raggedright
\strut
\end{minipage} & \begin{minipage}[t]{0.06\columnwidth}\raggedright
\strut
\end{minipage} & \begin{minipage}[t]{0.05\columnwidth}\raggedright
51,364\strut
\end{minipage} & \begin{minipage}[t]{0.04\columnwidth}\raggedright
44,091\strut
\end{minipage} & \begin{minipage}[t]{0.10\columnwidth}\raggedright
\strut
\end{minipage} & \begin{minipage}[t]{0.09\columnwidth}\raggedright
\strut
\end{minipage} & \begin{minipage}[t]{0.11\columnwidth}\raggedright
\strut
\end{minipage}\tabularnewline
\begin{minipage}[t]{0.03\columnwidth}\raggedright
3\strut
\end{minipage} & \begin{minipage}[t]{0.10\columnwidth}\raggedright
39--65\strut
\end{minipage} & \begin{minipage}[t]{0.07\columnwidth}\raggedright
6,485\strut
\end{minipage} & \begin{minipage}[t]{0.08\columnwidth}\raggedright
\strut
\end{minipage} & \begin{minipage}[t]{0.06\columnwidth}\raggedright
\strut
\end{minipage} & \begin{minipage}[t]{0.05\columnwidth}\raggedright
44,091\strut
\end{minipage} & \begin{minipage}[t]{0.04\columnwidth}\raggedright
36,818\strut
\end{minipage} & \begin{minipage}[t]{0.10\columnwidth}\raggedright
\strut
\end{minipage} & \begin{minipage}[t]{0.09\columnwidth}\raggedright
\strut
\end{minipage} & \begin{minipage}[t]{0.11\columnwidth}\raggedright
\strut
\end{minipage}\tabularnewline
\begin{minipage}[t]{0.03\columnwidth}\raggedright
4\strut
\end{minipage} & \begin{minipage}[t]{0.10\columnwidth}\raggedright
65--91\strut
\end{minipage} & \begin{minipage}[t]{0.07\columnwidth}\raggedright
19,230\strut
\end{minipage} & \begin{minipage}[t]{0.08\columnwidth}\raggedright
\strut
\end{minipage} & \begin{minipage}[t]{0.06\columnwidth}\raggedright
\strut
\end{minipage} & \begin{minipage}[t]{0.05\columnwidth}\raggedright
36,818\strut
\end{minipage} & \begin{minipage}[t]{0.04\columnwidth}\raggedright
29,546\strut
\end{minipage} & \begin{minipage}[t]{0.10\columnwidth}\raggedright
\strut
\end{minipage} & \begin{minipage}[t]{0.09\columnwidth}\raggedright
\strut
\end{minipage} & \begin{minipage}[t]{0.11\columnwidth}\raggedright
\strut
\end{minipage}\tabularnewline
\begin{minipage}[t]{0.03\columnwidth}\raggedright
5\strut
\end{minipage} & \begin{minipage}[t]{0.10\columnwidth}\raggedright
91--117\strut
\end{minipage} & \begin{minipage}[t]{0.07\columnwidth}\raggedright
37,878\strut
\end{minipage} & \begin{minipage}[t]{0.08\columnwidth}\raggedright
\strut
\end{minipage} & \begin{minipage}[t]{0.06\columnwidth}\raggedright
\strut
\end{minipage} & \begin{minipage}[t]{0.05\columnwidth}\raggedright
29,546\strut
\end{minipage} & \begin{minipage}[t]{0.04\columnwidth}\raggedright
22,273\strut
\end{minipage} & \begin{minipage}[t]{0.10\columnwidth}\raggedright
\strut
\end{minipage} & \begin{minipage}[t]{0.09\columnwidth}\raggedright
\strut
\end{minipage} & \begin{minipage}[t]{0.11\columnwidth}\raggedright
\strut
\end{minipage}\tabularnewline
\begin{minipage}[t]{0.03\columnwidth}\raggedright
6\strut
\end{minipage} & \begin{minipage}[t]{0.10\columnwidth}\raggedright
117--143\strut
\end{minipage} & \begin{minipage}[t]{0.07\columnwidth}\raggedright
63,427\strut
\end{minipage} & \begin{minipage}[t]{0.08\columnwidth}\raggedright
\strut
\end{minipage} & \begin{minipage}[t]{0.06\columnwidth}\raggedright
\strut
\end{minipage} & \begin{minipage}[t]{0.05\columnwidth}\raggedright
22,273\strut
\end{minipage} & \begin{minipage}[t]{0.04\columnwidth}\raggedright
15,000\strut
\end{minipage} & \begin{minipage}[t]{0.10\columnwidth}\raggedright
\strut
\end{minipage} & \begin{minipage}[t]{0.09\columnwidth}\raggedright
\strut
\end{minipage} & \begin{minipage}[t]{0.11\columnwidth}\raggedright
\strut
\end{minipage}\tabularnewline
\bottomrule
\end{longtable}

Total ideal Δv (altitude-varying Isp) = \_\_\_\_\_\_\_\_\_\_ ft/s

V\_bo = Total ideal Δv − 3540 ft/s = \_\_\_\_\_\_\_\_\_\_ ft/s

Compare to 1.0's constant-Isp result (V\_bo ≈ 6284 ft/s): difference =
\_\_\_\_\_\_\_\_\_\_ ft/s ( \_\_\_\_\_\_ \%)

\end{document}
